This calculation retains the arithmetic coefficient complex over the
absolute pointed base $\mathbf F_{1,\tau}=\{\tau,1\}$. It computes the
signed arithmetic/reference determinant term on the original finite
theta source. The new formulas are a complete first and second variation,
an exact four-window projection trace, and a quantitative bound in terms
of the actual common source Gram. No arithmetic source multiplier is
assumed bounded above or bounded away from zero.

\section{Original objects and the fixed comparison}

The retained tau-base construction has the four-point chart
$+,-<\eta<\sigma$ and the coefficient sheaf $\mathcal T$ with
$(\mathcal T_+,\mathcal T_-,\mathcal T_\eta,\mathcal T_\sigma)
=(V,V,\mathscr B,0)$. Its two nonzero restrictions are $\Theta$ and
$J\Theta=\Theta\mathcal F$. The actual projective resolution gives
\[
 \mathsf A_\tau(\mathcal T)
 =[V\oplus V\xrightarrow{\Theta(\phi-\widehat\psi)}\mathscr B],
 \qquad H^1\mathsf A_\tau(\mathcal T)=Q=\mathscr B/\Theta V.
 \tag{AT1}
\]
These are the fixed source objects of the original Tau Base note,
equations (18)--(25), supplied in the source closure. The tensor
resolution in its equations (46)--(48) retains the original differential
signs and identifies its top cohomology with $Q^{\otimes k}$.
The operation here uses that pushforward, with its arithmetic classes.
Its separately computed restriction to $\sigma$ remains an available
comparison; it is not substituted for (AT1).

Fix an original actual packet $h$ of full zero orders of $g=2\xi$,
and a positive integer $k$. Retain the full tensor-sum operator,
its original cyclic annihilator $\chi=\chi_{h,k}$ of degree $q\ge1$,
and $E=\C[S]/(\chi)$ in the increasing-power remainder frame.
The original construction supplies
\[
 \mathcal V P=P(D_1+\cdots+D_k)F_h^{\otimes k},\qquad
 J^{(k)}\mathcal V P=\eta_{h,k}J_NP,\qquad
 \eta_{h,k}[P]=\upsilon_h^{\otimes k}P(A_h^{(k)})1,
 \tag{AT2}
\]
where $\upsilon_h=j_h(g/h)$ is the full invertible Taylor unit,
$J_N$ is literal remainder modulo $\chi$, and
$\eta_{h,k}:E\hookrightarrow E_h^{\otimes k}$ is the original
injective full-jet observation, $E_h=\C[s]/(h)$.
The original section $\sigma_h:E_h\hookrightarrow Q$ gives the
separate arithmetic cohomology injection and equality
\[
 q^{(k)}\mathcal V P=\sigma_h^{\otimes k}\eta_{h,k}J_NP,
 \qquad \sigma_h^{\otimes k}\eta_{h,k}:E\hookrightarrow Q^{\otimes k}.
 \tag{AT2a}
\]
The supplied arithmetic-input and cyclic-sum proofs construct these
maps; all operations below keep them fixed. In
particular the sum coordinate is $S=k/2+iu$, with no division by $k$.

For $N\ge q-1$, set $\cP_N=\C[S]_{\le N}$,
$\cP_{-1}=0$, and let
$B_N:\cP_{N-q}\to\cP_N$ be multiplication by the literal monic
$\chi$. Write $s_N:E\to\cP_N$ for the degree-less-than-$q$
remainder section. Polynomial division proves
\[
 0\longrightarrow\cP_{N-q}\xrightarrow{B_N}\cP_N
       \xrightarrow{J_N}E\longrightarrow0,
 \qquad J_Ns_N=I_E.                         \tag{AT3}
\]
The arithmetic norm has density
$m_1(u)=m_{h,k}(u)=w_h^{*k}(u)$,
$w_h(u)=|(g/h)(1/2+iu)|^2/(2\pi)$.
The retained Gamma reference has density
\[
 m_0(u)=r_{1/4,k}(u)
 =\frac{(\sqrt{2\pi})^k}{c_{k/4}}
       \frac{|\Gamma(k/4+iu/2)|^2}{2\pi},
 \qquad c_a=2^{1-2a}\Gamma(2a).              \tag{AT4}
\]
Its mass is $(2\pi)^{k/2}$; the arithmetic mass is the original
$\mu_h^k$. The source proofs included with this calculation establish
finite moments of every degree and positivity almost everywhere for
both densities. The calculation below only integrates polynomials
against these two actual measures and their sum. It needs no estimate
on the pointwise ratio $m_1/m_0$.

For $0\le t\le1$ define, in the unchanged coefficient frame,
\[
 \begin{gathered}
 m_t=(1-t)m_0+tm_1,\quad \Delta m=m_1-m_0,
 \quad v_N(u)=(1,S(u),\ldots,S(u)^N),\\
 M_N(t)=\int v_N(u)^*v_N(u)m_t(u)\,du,
 \quad \Delta M_N=M_N(1)-M_N(0),\\
 H_N(t)=B_N^*M_N(t)B_N.                     \end{gathered}
 \tag{AT5}
\]
Every $M_N(t)$ is positive definite: a nonzero polynomial has finitely
many real zeros, and $m_t$ is positive almost everywhere. $B_N$ is
injective, so $H_N(t)$ is positive definite on its actual relation
domain. When $N=q-1$, that domain is zero; its determinant is one,
its unique endomorphism is its inverse, and all compositions through
that domain are zero. The real parameter $t$ interpolates norms on
(AT3); it changes neither the absolute base nor a cochain map.

\section{Canonical sections, determinant lines and full variations}

\begin{theorem}
On the fixed exact sequence (AT3), its least-norm section and quotient
Gram are
\[
 \begin{gathered}
 R_N(t)=s_N-B_NH_N(t)^{-1}B_N^*M_N(t)s_N,\\
 J_NR_N(t)=I_E,\quad B_N^*M_N(t)R_N(t)=0,\\
 G_N(t)=R_N(t)^*M_N(t)R_N(t)
       =(J_NM_N(t)^{-1}J_N^*)^{-1},\\
 V_N(t):=\det G_N(t)=\frac{\det M_N(t)}{\det H_N(t)}.
                                                        \end{gathered}
 \tag{AT6}
\]
In the original polynomial frame the exact first and second variations
are
\[
 \begin{gathered}
 \dot R_N=-B_NH_N^{-1}B_N^*\Delta M_NR_N,\qquad
 \dot G_N=R_N^*\Delta M_NR_N,\\
 \ddot G_N=-2R_N^*\Delta M_NB_NH_N^{-1}B_N^*\Delta M_NR_N
            \preceq0.                                \end{gathered}
 \tag{AT7}
\]
The exact logarithmic variation, with both trace terms retained, is
\[
 \begin{split}
 \frac{d}{dt}\log V_N
 &=\Tr(M_N^{-1}\Delta M_N)
   -\Tr(H_N^{-1}B_N^*\Delta M_NB_N)\\
 &=\Tr(G_N^{-1}R_N^*\Delta M_NR_N).             \end{split}
 \tag{AT8}
\]
Its second derivative is
\[
 \begin{split}
 \frac{d^2}{dt^2}\log V_N
 ={}&-2\Tr(G_N^{-1}R_N^*\Delta M_NB_NH_N^{-1}
                         B_N^*\Delta M_NR_N)\\
    &-\Tr(G_N^{-1}\dot G_NG_N^{-1}\dot G_N)\le0.
                                                        \end{split}
 \tag{AT9}
\]
\end{theorem}

\begin{proof}
Every lift of $x\in E$ is $s_Nx+B_Ny$. Expanding its squared norm
and completing its finite quadratic expression gives the unique
minimizer $y=-H_N^{-1}B_N^*M_Ns_N$. This proves the first two
lines of (AT6). For the inverse expression, every polynomial vector
$p$ decomposes as $B_Ny+R_NJ_Np$. Orthogonality therefore gives
$R_N^*M_Np=G_NJ_Np$ for every $p$. Taking adjoints yields
$M_NR_N=J_N^*G_N$. Substituting this into $J_NR_N=I_E$ gives
$(J_NM_N^{-1}J_N^*)G_N=I_E$, which proves the inverse expression.

The determinant calculation preserves orientation. Put
$r=N-q+1$ and $L_N=[B_N\ s_N]$. In the ordered monomial bases,
$[s_N\ B_N]$ has determinant one, since $\chi$ is monic and the
lower diagonal block is triangular with diagonal one. Thus
$\det L_N=(-1)^{qr}$. The upper-left block of $L_N^*M_NL_N$
is $H_N$, and its Schur complement is $G_N$. Consequently
$\det H_N\det G_N=|(-1)^{qr}|^2\det M_N=\det M_N$.
Equivalently the map
\[
 \begin{gathered}
 \det\cP_{N-q}\otimes\det E\longrightarrow\det\cP_N,\\
 \quad (y_1\wedge\cdots\wedge y_r)\otimes(x_1\wedge\cdots\wedge x_q)
 \longmapsto
 B_Ny_1\wedge\cdots\wedge B_Ny_r\wedge s_Nx_1\wedge\cdots\wedge s_Nx_q
 \end{gathered}
                                                        \tag{AT10}
\]
is the fixed determinant-line isomorphism. Replacing any lift by a
relation does not change its wedge. Orthogonal lifts give the product
metric just computed. This proves the determinant identity including
the empty relation space, without changing any scalar mass.

Differentiate $J_NR_N=I$; then $\dot R_N$ lies in $\im B_N$.
Differentiating $B_N^*M_NR_N=0$ determines its unique coefficient
and proves the first identity of (AT7). On differentiating
$R_N^*M_NR_N$, the two terms involving $\dot R_N$ vanish by
that orthogonality. A second differentiation gives the two equal
negative terms in (AT7), since $\Delta M_N$ is Hermitian. This
proves the claimed second derivative and its sign. Applying the
finite determinant derivative to (AT6) gives the first equality
of (AT8); applying it to $G_N$ gives the second.
One more derivative gives (AT9). The first trace there is the squared
Hilbert--Schmidt norm of
$H_N^{-1/2}B_N^*\Delta M_NR_NG_N^{-1/2}$; the second is the
squared Hilbert--Schmidt norm of
$G_N^{-1/2}\dot G_NG_N^{-1/2}$. These identities establish their
signs and retain the original forms. All inverse matrices are smooth
on a neighborhood of the compact interval $[0,1]$, by positivity
and continuity. No infinite-dimensional inverse has been used.
\end{proof}

The derivative $\dot R_N$ in (AT7) is a literal polynomial relation.
Applying the original $\mathcal V$ in (AT2) therefore retains its
original theta primitive and sends its arithmetic class to zero;
$J^{(k)}\mathcal V R_N(t)=\eta_{h,k}$ for every $t$.
At each nonempty original support label $A$, the cochain identities
lift by $(A,P)\mapsto(A,fP)$, while $\tau\mapsto\tau$.
Thus a relation-valued derivative has a supported zero class at
label $A$. Its source vector and its full nonnegative energy in
(AT7)--(AT9) remain present. These formulas compute their
contribution rather than deleting it after taking cohomology.

For every independent coefficient mask $B_c\subseteq I_n$, retain
the separate outer source and leg label $(W,S)$ and apply each of
these linear maps in every active coordinate. The supported-zero
fibre of $f[B_c]$ is precisely $(\ker f)^{B_c}$ at that same label.
For $B_c\subseteq C_c$, zero insertion commutes with every such
map: both composites use $f(x_i)$ in an old coordinate and $f(0)=0$
in a new coordinate. It is a labelled face arrow $e_{B_c}\mapsto
e_{C_c}$, not a collapse of the two faces. At empty coefficient
support the zero amplitude space retains $(W,S;\varnothing,0)$,
distinct from global external $\tau_{\rm out}$.
In each active coordinate the relation-valued derivative
is the original $\dot R_N$ and has the same theta primitive.
The equation $(J^{(k)}\mathcal V R_N)[B_c]=\eta_{h,k}[B_c]$
holds because it is the already proved equality in every coordinate.
The source derivative and the forms (AT7)--(AT9) remain the original
ones on those vectors.

\section{The signed four-endpoint term on one retained source}

Let $N_*=2q$. Every space in (AT3) for the four endpoints sits in
$\cP_{N_*}$ by its actual coefficient inclusion. Give this common
space the form $M_*(t)=M_{N_*}(t)$. Write $P_N(t)$ for its
orthogonal projection onto $\cP_N$, and $Q_N(t)$ for its
orthogonal projection onto the relation subspace
$\cR_N=\chi\cP_{N-q}$. The latter is zero at $N=q-1$.
For the coefficient inclusion $I_N:\cP_N\hookrightarrow\cP_{N_*}$,
the formulas, with their domains retained, are
\[
 \begin{gathered}
 P_N=I_NM_N^{-1}I_N^*M_*,\qquad
 Q_N=I_NB_NH_N^{-1}B_N^*I_N^*M_*,\qquad
 C(t)=M_*(t)^{-1}\Delta M_*.                   \end{gathered}
 \tag{AT11}
\]
Both $P_N$ and $Q_N$ are $M_*(t)$-self-adjoint idempotents with
their displayed images, verified by substitution. Since
$\cR_N\subset\cP_N$, $P_NQ_N=Q_NP_N=Q_N$. Their difference is
the orthogonal projection onto the finite canonical quotient
section and has trace $q$. Cyclicity of finite matrix trace in
(AT11) proves the exact common-source version of (AT8):
\[
 \frac{d}{dt}\log V_N(t)=\Tr((P_N(t)-Q_N(t))C(t)).
                                                        \tag{AT12}
\]

Put $i_0=q-1$, $i_1=q$, $j_0=2q-1$, and $j_1=2q$.
Define the two positive operators
\[
 \begin{split}
 U(t)&=(Q_{j_0}-Q_{i_0})+(Q_{j_1}-Q_{i_1}),\\
 W(t)&=(P_{j_0}-P_{i_0})+(P_{j_1}-P_{i_1}).    \end{split}
 \tag{AT13}
\]
Each difference in this display is an orthogonal projection,
because the two spaces in its pair are nested in the same
Hilbert space. Each has rank $q$. The sums need not be
projections and are not asserted to be orthogonal sums; they
are positive operators with $\Tr U=\Tr W=2q$.

Write
\[
 \begin{gathered}
 \mathcal B(t)=\log\frac{V_{i_0}(t)V_{i_1}(t)}{V_{j_0}(t)V_{j_1}(t)},
 \quad T_N=V_N(1)/V_N(0),\\
 \Delta\mathcal B=\mathcal B(1)-\mathcal B(0)
 =\log\frac{T_{q-1}T_q}{T_{2q-1}T_{2q}}.       \end{gathered}
 \tag{AT14}
\]
Subtracting the four exact formulas (AT12), with their original
signs, and integrating gives
\[
 \boxed{\displaystyle
 \Delta\mathcal B=\int_0^1\Tr((U(t)-W(t))C(t))\,dt .}
                                                        \tag{AT15}
\]
This is the actual arithmetic correction in the audit's formula
(54). It compares the two added relation windows with the two
added polynomial windows. Concavity of each individual
$\log V_N(t)$ in (AT9) does not assign a sign to their signed
combination in (AT15); that combination has now been computed.

There is also an exact pointwise integral, using no division by
the arithmetic density. Define
\[
 \begin{gathered}
 k_N(t,u)=v_N(u)M_N(t)^{-1}v_N(u)^*,\\
 r_N(t,u)=|\chi(S(u))|^2
 v_{N-q}(u)H_N(t)^{-1}v_{N-q}(u)^*,\\
 d_N(t,u)=k_N(t,u)-r_N(t,u),\qquad r_{q-1}(t,u)=0,\\
 d_N(t,u)=v_N(u)R_N(t)G_N(t)^{-1}R_N(t)^*v_N(u)^*\ge0.
                                                        \end{gathered}
 \tag{AT16}
\]
The last equality follows by writing the orthogonal splitting
of $M_N^{-1}$ as
$B_NH_N^{-1}B_N^*+R_NG_N^{-1}R_N^*$; multiplying by $M_N$
and evaluating on $\im B_N\oplus\im R_N$ proves that inverse
identity. In particular $\int d_N(t,u)m_t(u)\,du=q$.
Consequently (AT8) also reads
\[
 \frac{d}{dt}\log V_N(t)=\int d_N(t,u)\Delta m(u)\,du,
\quad
 \Delta\mathcal B=\int_0^1\int
 (d_{i_0}+d_{i_1}-d_{j_0}-d_{j_1})(t,u)\Delta m(u)\,du\,dt.
                                                        \tag{AT17}
\]
All coefficients of these fixed-degree polynomials in $u$ and
$\bar S(u)$ are bounded for $t\in[0,1]$. Since
$|\Delta m|\le m_0+m_1$ and these measures have the required
finite moments, absolute integrability proves the exchange of
integrals. This remains valid at every zero of $m_1$.

\section{An exact finite-source bound for the correction}

Let $b_1\le\cdots\le b_{2q+1}$ be every positive eigenvalue of
$D_*=M_*(0)^{-1}M_*(1)$ in the original $M_*(0)$ metric, and put
$\kappa_*=b_{2q+1}/b_1$. These data are on the unchanged common
source $H=\cP_{2q}$; neither packet nor tensor order is suppressed.
For the exact $U,W,C$ in (AT11)--(AT15), let $c_j(t)$ be the
ordered eigenvalues of $C(t)$, $d(t)=c_{2q+1}(t)-c_1(t)$, and
$s(t)=\dim(\operatorname{ran}U\cap\operatorname{ran}W)$. All norms
and projections in the following formulas use $M_*(t)$:
\[
\begin{aligned}
S_{\rm spec}(t)&=c_{q+2}(t)-c_q(t)
 +2\sum_{j=1}^{q-1}(c_{2q+2-j}(t)-c_j(t)),\\
S_{\rm tr}(t)&=\tfrac12d(t)\|U(t)-W(t)\|_{1,t},\\
S_{\rm ov}(t)&=\tfrac12d(t)\min\{4q-2s(t),
 \sqrt{(2q+1)(8q-4-2\Tr(U(t)W(t)))}\},\\
\mathfrak C_*&=\int_0^1\min\{S_{\rm spec}(t),S_{\rm tr}(t)\}\,dt,\\
\mathfrak C_*^{\rm ov}&=\int_0^1\min\{S_{\rm spec}(t),S_{\rm ov}(t)\}\,dt,\\
L_q(D_*)&=\log\frac{b_{q+2}}{b_q}
 +2\sum_{j=1}^{q-1}\log\frac{b_{2q+2-j}}{b_j}.
\end{aligned}\tag{AT18a}
\]
The trace norm is the sum of absolute eigenvalues of the displayed
self-adjoint endomorphism. For $q=1$ the sums are empty. Under
$x=t$, these are literally the ACM/SP operators: their relation
matrix is $I_NB_N$, hence their projector is exactly (AT11);
their density $r_{1/4}^{*k}$ equals (AT4), with
$c_{1/4}=\sqrt{2\pi}$ and its full mass. Thus
$\mathfrak C_*^{\rm ov}$ is the original ACM pointwise minimum.
The sharper $\mathfrak C_*$ also retains SP's already proved exact
trace norm; no new source form is chosen.

\begin{theorem}
The actual signed correction satisfies
\[
\boxed{|\Delta\mathcal B|\le J_*\le\mathfrak C_*
 \le\mathfrak C_*^{\rm ov}\le L_q(D_*)
 \le(2q-1)\log\kappa_* .}\tag{AT18}
\]
Here the full actual-moment proof of the additional $J_*$ entry is
given in (AT21a)--(AT21e). Consequently its original endpoint quantities obey
\[
\mathcal B^\Ga_{h,k}-J_*
 \le\mathcal B^\ar_{h,k}
 \le\mathcal B^\Ga_{h,k}+J_* .\tag{AT19}
\]
\end{theorem}
\begin{proof}
Multiplication by the same monic $\chi$ gives the nested relation
flag $\chi\cP_0\subset\chi\cP_{q-1}\subset\chi\cP_q$ of
dimensions $1,q,q+1$. On its orthogonal increments $U$ acts by
$1,2,1$, and acts by zero on the $q$-dimensional complement.
The polynomial flag $\cP_{q-1}\subset\cP_q\subset\cP_{2q-1}
\subset\cP_{2q}$ gives the same full spectrum for $W$:
$0^q,1^2,2^{q-1}$. Repeated spaces when $q=1$ give zero-dimensional
increments, not omitted endpoint maps. In particular
\[
\Tr U=\Tr W=2q,\quad \Tr U^2=\Tr W^2=4q-2,
\quad\Tr((U-W)^2)=8q-4-2\Tr(UW).
\]
For a rank-$r$ orthogonal projection $P$, its diagonal entries in
an orthonormal eigenbasis of $C$ lie in $[0,1]$ and sum to $r$.
Moving this fixed nonnegative mass to the $r$ smallest or largest
eigenvalues proves
$\sum_{j=1}^r c_j\le\Tr(PC)\le\sum_{j=2q+2-r}^{2q+1}c_j$.
Each of $U,W$ is its range projection of rank $q+1$ plus its
eigenvalue-two projection of rank $q-1$. Applying these two bounds
and subtracting the complete sums gives $|\Tr((U-W)C)|\le S_{\rm spec}$.

Since $\Tr(U-W)=0$, subtracting the midpoint of the extremes of
$C$ leaves the trace pairing unchanged. In an orthonormal eigenbasis
of $U-W$ every resulting diagonal entry of $C$ has absolute value
at most $d/2$. Therefore $|\Tr((U-W)C)|\le S_{\rm tr}$.
The two ranges have dimension $q+1$ in dimension $2q+1$, so $s\ge1$.
Let $E$ be their intersection projection. The inequalities
$U\succeq E$, $W\succeq E$ follow from their nonzero eigenvalues
being at least one. The two positive differences have trace $2q-s$;
their difference is $U-W$. Thus $\|U-W\|_1\le4q-2s$.
Cauchy--Schwarz on its $2q+1$ real eigenvalues proves the second
trace-norm upper bound in $S_{\rm ov}$. We have proved pointwise
\[
|\Tr((U-W)C)|\le\min\{S_{\rm spec},S_{\rm tr}\}
 \le\min\{S_{\rm spec},S_{\rm ov}\}.\tag{AT20}
\]
All the matrix entries are smooth on the compact interval. Eigenvalues
and trace norms are continuous. The intersection dimension is Borel,
since rank strata are specified by vanishing and nonvanishing minors.
The bounded nonnegative functions above are consequently integrable.

In the fixed coefficient space,
$C(t)=((1-t)I+tD_*)^{-1}(D_*-I)$, so
\[
c_j(t)=\frac{b_j-1}{1-t+tb_j},\qquad
\int_0^1c_j(t)\,dt=\log b_j,\qquad
\int_0^1d(t)\,dt=\log\kappa_* .\tag{AT21}
\]
The derivative of the fraction with respect to $b_j$ is
$(1-t+tb_j)^{-2}>0$, proving its order through both endpoints.
Integrating (AT20) in (AT15) proves the first inequalities of
(AT18). Integrating $S_{\rm spec}$ gives $L_q(D_*)$ exactly.
Its logarithmic ratios are each between zero and $\log\kappa_*$;
their total multiplicity is $1+2(q-1)=2q-1$. This proves the last
inequality. Substitution of (AT14) gives the additive endpoint consequence;
the stronger $J_*$ version of (AT19) is proved in (AT21a)--(AT21e).
The preceding coarser $2q\log\kappa_*$ estimate is a corollary
because $\log\kappa_*\ge0$; its weaker coefficient is not used
in the current endpoint consequences.
\end{proof}

\subsection{The same control in the original nonlinear endpoint coordinate}

The additional moment structure is retained in this step. The two
endpoint forms are the original positive measures on $S=k/2+iy$;
the proof also applies to a specified positive atomic measure on
that line when its degree-$2q$ moment Gram is positive definite.
No arbitrary coefficient form is assigned this moment structure.
Write $a=2q-1$, $B(t)=\mathcal B(t)$, and
$f_a(B)=a\sqrt{1-e^{-B/a}}$ for $B>0$. The exact additional
functions are
\[
\begin{aligned}
S_{\rm sin}(t)&=d(t)\sum_{\ell=1}^a\sqrt{1-e^{-\lambda_\ell(t)}},
\qquad S_{\rm ang}(t)=f_a(B(t))d(t),\\
\Sigma(t)&=\Tr(C(t)^2)-\frac{(\Tr C(t))^2}{2q+1},\\
S_{\rm HS}(t)&=\sqrt{\Sigma(t)\Tr((U(t)-W(t))^2)},\\
J_*&=\int_0^1\min\{S_{\rm spec},S_{\rm tr},S_{\rm ov},
                      S_{\rm HS},S_{\rm sin},S_{\rm ang}\}\,dt,\\
I_*&=\int_0^1\min\{d(t),S_{\rm spec}/f_a(B(t)),
                    S_{\rm tr}/f_a(B(t)),\\
&\hspace{22mm}S_{\rm ov}/f_a(B(t)),S_{\rm HS}/f_a(B(t)),S_{\rm sin}/f_a(B(t))\}\,dt,\\
H_a(B)&=2\operatorname{arcosh}(e^{B/(2a)}),\qquad
\Psi_a(z)=2a\log\cosh(z/2)\quad(z\ge0).
\end{aligned}\tag{AT21a}
\]
The $\lambda_\ell$ are the complete crossed-pair losses specified
below, with the forced zero contribution separately retained.
The additional centered Hilbert--Schmidt estimate is proved on the
same original $M_*(t)$ metric as follows. Since $\Tr(U-W)=0$,
subtract $\Tr(C)I/(2q+1)$ from $C$ inside the trace pairing.
Its squared Hilbert--Schmidt norm is exactly $\Sigma(t)\ge0$.
For Hermitian matrices in an $M_*(t)$-orthonormal basis,
$\Tr(AB)=\sum_{i,j}A_{ij}\overline{B_{ij}}$; Cauchy--Schwarz
on those finitely many entries proves
$|\Tr(AB)|\le\sqrt{\Tr(A^2)\Tr(B^2)}$. Applying this to the
centered $C$ and $U-W$ proves $|B'(t)|\le S_{\rm HS}(t)$.
The metric is retained on both operators, so this proof changes
neither the source nor its arithmetic observation.
We prove the complete moment and projection facts used here. For
cutoffs $i<j$, regard the original minimum sections $R_i,R_j$ as
maps to $H$. Their difference lies in $\ker J_j$, and $R_j$ is
orthogonal to that kernel, so
\[
R_i^*MR_j=G_j,\qquad
G_i-G_j=(R_i-R_j)^*M(R_i-R_j)\succeq0.
\]
The orthonormal frames $X_i=R_iG_i^{-1/2}$,
$X_j=R_jG_j^{-1/2}$ have overlap
$X_i^*MX_j=G_i^{-1/2}G_j^{1/2}$. Its singular values
$\sigma_\ell$ are positive and at most one, and
$\prod_\ell\sigma_\ell^2=\det G_j/\det G_i$.
To compute the projection difference exactly, choose paired unit
vectors with overlap $\sigma_\ell$. When $\sigma_\ell<1$, the
unit vector $(y_\ell-\sigma_\ell x_\ell)/\sqrt{1-\sigma_\ell^2}$
is orthogonal to the first space. In this two-dimensional plane the
two projection matrices differ by a trace-zero matrix of determinant
$-(1-\sigma_\ell^2)$, with eigenvalues
$\pm\sqrt{1-\sigma_\ell^2}$. An overlap one gives the identical
direction and zero difference. Orthogonal complements give further
zeros. Thus its trace norm is exactly
$2\sum_\ell\sqrt{1-\sigma_\ell^2}$, without changing the source norm.

Apply this to the crossed pairs $(q-1,2q)$ and $(q,2q-1)$.
Their differences sum to $U-W$ by (AT11)--(AT13). The second
pair's two $q$-dimensional section spaces lie in
$\cP_{2q-1}\cap(\chi\cP_0)^\perp$, of dimension $2q-1$;
therefore they share a direction. Retain its overlap one in the
full second list. The remaining $a=2q-1$ logarithmic losses
$\lambda_\ell=-\log\sigma_\ell^2$ include every further zero,
and their sum is $B(t)$. The trace-norm triangle inequality gives
\[
S_{\rm tr}(t)\le d(t)\sum_{\ell=1}^{a}\sqrt{1-e^{-\lambda_\ell}}
 \le f_a(B(t))d(t)=S_{\rm ang}(t).\tag{AT21b}
\]
The last inequality follows from concavity of
$z\mapsto\sqrt{1-e^{-z}}$: its second derivative for $z>0$ is
$-e^{-z}/(2\sqrt{1-e^{-z}})-e^{-2z}/(4(1-e^{-z})^{3/2})<0$,
and continuity includes zero. Jensen's inequality on the full
$a$-entry list proves the displayed bound.

The number $B(t)$ is strictly positive. If it were zero, both
nonnegative logarithmic determinant ratios would vanish. All relative
Gram eigenvalues are at least one, so each ratio one forces
$G_i=G_j$, hence $R_i=R_j$ by the squared norm above. In particular
$1\in\operatorname{im}R_{q-1}=\operatorname{im}R_{2q}$ would be
orthogonal to $\chi\cP_q$. For
$\chi(S)=\sum a_jS^j$ let
$\chi^{\#_k}(S)=\sum\overline{a_j}(k-S)^j$. On the actual line
this is $\overline{\chi(S)}$. The asserted orthogonality would give
$0=\int|\chi(S)|^2d\mu_t$, contradicting the positive moment
Gram of the nonzero polynomial $\chi$. This proof also covers the
specified positive atomic source. Smoothness and compactness now
give a positive minimum for $B(t)$, so all denominators in (AT21a)
are bounded away from zero on this fixed path.

Since $S_{\rm tr}\le S_{\rm ov}$ and (AT21b) holds, pointwise
comparison proves
\[
J_* =\int_0^1\min\{S_{\rm spec},S_{\rm tr},S_{\rm HS}\}\,dt
 \le\mathfrak C_*,\qquad
I_* =\int_0^1\frac{\min\{S_{\rm spec},S_{\rm tr},S_{\rm HS}\}}{f_a(B(t))}\,dt
 \le\int_0^1d(t)\,dt=\log\kappa_* .\tag{AT21c}
\]
Differentiation gives $H_a'(B)=1/f_a(B)>0$; substitution gives
$\Psi_a(H_a(B))=B$ and $H_a(\Psi_a(z))=z$ for $z>0$.
Intersecting the original (AT20) estimate with the proved
$S_{\rm HS}$ bound gives $|B'|\le\min(S_{\rm spec},S_{\rm tr},S_{\rm HS})$.
It proves $|B(1)-B(0)|\le J_*$. Applying the exact derivative
to that same pointwise bound and integrating proves
\[
|H_a(B(1))-H_a(B(0))|\le I_* .\tag{AT21d}
\]
This proves the recursive nonlinear refinement from the original
source maps rather than assuming a stronger endpoint bound.
Put $H_0=H_a(B(0))$ and define the actual current endpoints
\[
\begin{aligned}
\mathfrak b_*^-&=\max\{0,B(0)-J_*,
                      \Psi_a(\max\{0,H_0-I_*\})\},\\
\mathfrak b_*^+&=\min\{B(0)+J_*,\Psi_a(H_0+I_*)\},\\
\mathfrak d_*^-&=\mathfrak b_*^--B(0),\qquad
\mathfrak d_*^+=\mathfrak b_*^+-B(0).
\end{aligned}\tag{AT21e}
\]
Monotonicity of $\Psi_a$ on $[0,\infty)$, the proved additive
$J_*$ bound, and (AT21d)
prove $\mathfrak b_*^-\le B(1)\le\mathfrak b_*^+$ and
$\mathfrak d_*^-\le\Delta\mathcal B\le\mathfrak d_*^+$.
These endpoints refine the earlier additive interval while retaining
every mass, angle and spectrum. The case $q=1$ has $a=1$, one
angle, and the displayed zero second-pair angle. No estimate
uniform in $k$ follows from these fixed-path integrals.

When $M_*(1)=cM_*(0)$ with its literal $c>0$, the exact source
form at time $t$ is $(1-t+tc)M_*(0)$. Each quotient determinant
is $(1-t+tc)^qV_N(0)$, so all four factors in (AT14) cancel
and $\Delta\mathcal B=0$. This is also forced by
$\kappa_*=1$ in (AT18). The individual masses and volume
factors remain as written; cancellation occurs only in this
specified ratio. For an arbitrary fixed source-coordinate change, write old
common-source coordinates as $T$ times new coordinates, and include
each original section in that source as $\mathsf R_N=I_NR_N$.
With the quotient frame fixed, the exact transported maps and forms are
\[
 \begin{gathered}
 M_*'=T^*M_*T,\qquad J_N'=J_NT,\\
 \mathsf R_N'=T^{-1}\mathsf R_N,
 \qquad G_N'=G_N,\qquad V_N'=V_N.
 \end{gathered}
\]
Here each endpoint source is its transported subspace of the
common source, and $J_NT$ denotes the restriction to that subspace.
The transported section remains the minimum section because the
change is an isometry between the two displayed forms and preserves
each affine remainder fibre. Direct substitution then gives
$(\mathsf R_N')^*M_*'\mathsf R_N'=\mathsf R_N^*M_*\mathsf R_N=G_N$.
If, in addition,
old quotient coordinates equal $A$ times new quotient coordinates,
then $J_N'=A^{-1}J_NT$, $\mathsf R_N'=T^{-1}\mathsf R_NA$, and
$G_N'=A^*G_NA$, $V_N'=|\det A|^2V_N$.
The same quotient-frame multiplier at all four endpoints cancels
in (AT14). The source change conjugates $D_*$ and $C$, and sends
each projection to $T^{-1}PT$; its trace and spectrum therefore
remain unchanged. These exact maps preserve (AT15)--(AT21).

For the reflection-stable quartet calculation of the supplied
Gamma audit, let its actual displacement be $\delta>0$ and
retain $q=[1+k(m-1)](k+1)^2$. Its fully written Gamma theorem
gives $\mathcal B^\Ga_{h,k}\ge
4q\log(\delta k/(2\sqrt5))$. Substitution in (AT19) gives the
new quantitative comparison on the same source:
\[
\begin{aligned}
 \mathcal B^\ar_{h,k}&\ge
 4q\log\!\left(\frac{\delta k}{2\sqrt5}\right)-J_*,\\
 \mathfrak C_*&\ge J_*\ge\max\left\{0,
 4q\log\!\left(\frac{\delta k}{2\sqrt5}\right)-\mathcal B^\ar_{h,k}\right\},\\
 \log\kappa_*&\ge\frac{1}{2q-1}\max\left\{0,
 4q\log\!\left(\frac{\delta k}{2\sqrt5}\right)-\mathcal B^\ar_{h,k}\right\}.
\end{aligned}\tag{AT22}
\]
The same original Gamma lower bound also propagates through the
nonlinear endpoint coordinate. Set
$L_\Gamma=\max\{0,4q\log(\delta k/(2\sqrt5))\}$.
The full original Gamma moment argument gives $B(0)\ge L_\Gamma$,
and (AT21e) then implies
\[
\mathcal B^{\ar}_{h,k}\ge
\max\{0,L_\Gamma-J_*,
 \Psi_a(\max\{0,H_a(L_\Gamma)-I_*\})\},\qquad a=2q-1.
\tag{AT22a}
\]
Here $H_a(0)=0$ is its continuous extension. Each function of the
initial value in (AT21e) is increasing, so substituting the proved
lower bound proves every displayed entry without replacing the
actual path integral $I_*$. The simpler first-line estimate (AT22)
is retained as its weaker corollary.

The quartet and actual packet restriction are retained in this
application; a reference test polynomial has not been declared
to divide $2\xi$. Equations (AT15) and (AT17) give the full
signed quantity behind this bound. In particular they provide
the exact integrand that an arithmetic cancellation calculation
must evaluate, with both relation windows and both polynomial
windows present. No subthreshold arithmetic bound or weight
purity is asserted by this source comparison alone.

\section{Source dependencies and precise status}

The complete retained Tau Base note supplies (AT1), its tensor
pushforward, right adjoint and supported maps. The original
arithmetic-input and cyclic-sum proofs supply (AT2), full packet
units and the theta primitives. The Gamma/source-density notes
supply the actual measures and finite moments in (AT4)--(AT5).
The new Deligne-bound note supplies the full quartet Gamma
lower bound used only in (AT22). Each is preserved in the
accompanying source closure, with its actual edition and hash.

The new proved calculations here are (AT6)--(AT21) on those
unchanged input objects, and their substitution (AT22). They
give an analytic calculation on the canonical finite complex
mapping to $\mathsf A_{\tau,k}(\mathcal T^{\boxtimes k})$.
They neither replace that pushforward by a generic stalk nor
deduce a verdict about the entire tau-based program from an
auxiliary model.
